\question[5]
Erkläre anhand eines eigenen Python-Beispiels den Unterschied zwischen den beiden
Möglichkeiten, eine Liste mit einer Schleife zu durchlaufen:

\begin{solutionbox}{8cm}
\begin{lstlisting}
# Laufvariable i durchläuft die Indizes:
a = [3, 4, 7]
for i in range(len(a)):
    print(i, a[i])

# Laufvariable x durchläuft die Elemente:
a = [3, 4, 7]
for x in a:
    print(x)
\end{lstlisting}

Unterschied: Bei der Index-Variante steht zusätzlich der
Index zur Verfügung (z.B. um a[i] zu verändern oder um auf Nachbarn zuzugreifen), bei der
Element-Variante erhält man direkt die Werte, der Code
ist kürzer und übersichtlicher.

\end{solutionbox}